\documentclass[12pt,a4paper]{article}
\usepackage[T1]{fontenc}
\usepackage[utf8]{inputenc}
\usepackage[brazil]{babel}
\usepackage{mathptmx}
\usepackage[a4paper,top=3cm,bottom=2cm,left=3cm,right=2cm]{geometry}
\usepackage{booktabs,tabularx,array,graphicx,caption,xcolor,enumitem,amsmath}
\usepackage{tikz}
\usetikzlibrary{positioning,arrows.meta,fit,automata,calc}
\usepackage[hidelinks]{hyperref}
\usepackage{listings}

\lstset{basicstyle=\ttfamily\footnotesize,columns=fullflexible,frame=single,breaklines=true,
  framesep=4pt,xleftmargin=6pt,xrightmargin=6pt}
\captionsetup{font=small,labelfont=bf}
\setlength{\parskip}{4pt}
\setlist{itemsep=2pt,topsep=3pt}
\newcolumntype{L}{>{\raggedright\arraybackslash}X}
\renewcommand{\arraystretch}{1.15}

\title{\textbf{Monitoramento ubíquo de frequência cardíaca com mudança de contexto pelo acelerômetro do smartphone}\\[4pt]
\large Um protótipo com ESP32 simulado, fila FIFO, banco temporal e aplicativo Android}
\author{Felipe Alves \and Felipe O. Carvalho \and Matheus Augusto \and Murilo Bernardo\\[4pt]
\small Instituto de Informática -- Universidade Federal de Goiás (UFG)\\
\footnotesize\texttt{felipe\_leao@discente.ufg.br}, \texttt{felipecarvalho@discente.ufg.br},\\
\footnotesize\texttt{matheus.medeiros@renapsi.org.br}, \texttt{m.bernardo@discente.ufg.br}\\[2pt]
\small Software para Sistemas Ubíquos (INF0483) -- 2026/2}
\date{27 de novembro de 2026 -- versão do Marco 4}

\begin{document}
\maketitle

\begin{abstract}
\noindent
Alertas de frequência cardíaca (FC) elevada perdem valor quando disparam em situações em que a FC alta é esperada, como durante o esforço físico.
Este trabalho apresenta um protótipo de monitoramento ubíquo em que a decisão sobre o estado de um paciente fictício combina dois produtores de eventos.
O primeiro é um ESP32 simulado no Wokwi, que mede FC sintética e aceleração.
O segundo é o acelerômetro do smartphone do cuidador, que informa se o paciente está em movimento.
Uma API em FastAPI recebe os dois fluxos em uma única fila FIFO, aplica uma regra com janela, cobertura mínima, persistência, histerese e validade de contexto, grava leituras, transições e comandos no TimescaleDB e devolve a resposta ao cuidador, que pode reconhecer o alerta e silenciar o atuador do dispositivo.
A autenticação usa HMAC-SHA256 no dispositivo e, no aplicativo, um token assinado obtido por login, sem chaves no APK, sobre HTTPS validado pelo Android.
A avaliação combina 49 testes automatizados, executados em integração contínua, com medições reprodutíveis.
O processamento de um lote leva cerca de 4,5\,ms dentro da API.
Com a FC em 120\,bpm por 10 minutos, o sistema gera um alerta sem o contexto do celular e nenhum quando o celular indica movimento.
A avaliação também revelou e permitiu corrigir um defeito na combinação dos dois produtores.
\end{abstract}

\noindent\textbf{Palavras-chave:} computação ubíqua; consciência de contexto; Internet das Coisas; monitoramento de saúde; processamento de eventos.

% ------------------------------------------------------------------
\section{Introdução}

A computação ubíqua propõe sistemas que se integram ao ambiente e ao cotidiano das pessoas, de modo que a tecnologia ``desapareça'' do primeiro plano \cite{weiser1991}.
Em saúde, essa visão se traduz em dispositivos vestíveis e sensores conectados que acompanham sinais do paciente continuamente e avisam alguém quando algo foge do esperado \cite{islam2015}.
O valor desse acompanhamento depende menos da quantidade de dados coletados do que da qualidade das decisões tomadas sobre eles.
Um sistema que alerta demais acaba ignorado: o fenômeno da \emph{fadiga de alarmes} é documentado em ambientes clínicos, nos quais a maior parte dos alarmes não exige ação \cite{sendelbach2013}.

A frequência cardíaca ilustra bem o problema.
Uma FC de 120\,bpm em repouso pode indicar uma situação que merece atenção, mas a mesma FC durante uma caminhada rápida é normal.
Um alerta baseado só no limiar de FC, portanto, dispara em situações irrelevantes.
O que distingue os dois casos é o \emph{contexto}, entendido como qualquer informação que caracteriza a situação de uma entidade relevante para a interação \cite{dey2001}.
Aqui, o contexto é se o paciente está ou não em movimento.

Este trabalho descreve o projeto ``Internet dos Coisos'', desenvolvido ao longo de quatro marcos da disciplina Software para Sistemas Ubíquos.
O sistema simula uma roupa inteligente com um ESP32 que mede FC e aceleração, envia as leituras a um servidor e acende um LED e um buzzer quando o servidor decide que o paciente precisa de atenção.
O cuidador acompanha o estado por um painel web e por um aplicativo Android e pode reconhecer o alerta, o que silencia o buzzer.
No quarto marco, o aplicativo passou a ser também o \emph{segundo produtor} de eventos: o acelerômetro do celular informa o movimento do paciente, e esse contexto suprime ou encerra o alerta de FC alta.

As contribuições deste trabalho são:
\begin{enumerate}
  \item a arquitetura de um sistema com dois produtores heterogêneos, um dispositivo embarcado e um smartphone, consumidos por uma única fila FIFO e por uma regra com estado;
  \item uma regra de decisão que combina janela deslizante, cobertura mínima, persistência, histerese, tempo mínimo por estado e um contexto externo com prazo de validade;
  \item um esquema de autenticação e autorização em que dispositivo, painel e aplicativo têm identidades e permissões distintas, sem segredos no repositório nem no APK;
  \item uma avaliação reprodutível, com testes automatizados em integração contínua e medições com relógio simulado, que revelou e permitiu corrigir um defeito na interação entre os dois produtores.
\end{enumerate}

O restante do texto está organizado assim.
A Seção~\ref{sec:fund} apresenta a fundamentação e os trabalhos relacionados.
A Seção~\ref{sec:arq} descreve a proposta e a arquitetura, e a Seção~\ref{sec:impl}, a implementação.
A Seção~\ref{sec:metodo} define a metodologia de avaliação, cujos resultados são discutidos na Seção~\ref{sec:result}.
A Seção~\ref{sec:concl} conclui o trabalho.

% ------------------------------------------------------------------
\section{Fundamentação e trabalhos relacionados}\label{sec:fund}

\subsection{Computação ubíqua e consciência de contexto}

Weiser descreveu a computação ubíqua como computadores embutidos no ambiente, usados sem atenção consciente \cite{weiser1991}.
Satyanarayanan apontou os desafios práticos dessa visão, entre eles a percepção do contexto do usuário, o uso inteligente de espaços e a adaptação a recursos variáveis \cite{satya2001}.
Dey define contexto como a informação usada para caracterizar a situação de uma entidade e considera um sistema \emph{consciente de contexto} quando usa essa informação para oferecer serviços relevantes à tarefa do usuário \cite{dey2001}.
Abowd et al.\ destacam quatro tipos primários de contexto: localização, identidade, tempo e atividade \cite{abowd1999}.
O protótipo usa dois deles: o \emph{tempo}, nas janelas, na persistência e na validade das leituras, e a \emph{atividade}, no movimento do paciente.

\subsection{Internet das Coisas em saúde e reconhecimento de atividade}

Islam et al.\ revisam arquiteturas de IoT para saúde e organizam essas soluções em sensores no corpo, gateways e serviços na nuvem, destacando requisitos de segurança e privacidade \cite{islam2015}.
O reconhecimento de atividade humana com sensores vestíveis, em especial acelerômetros, é uma área consolidada.
Lara e Labrador revisam técnicas que vão de limiares sobre a magnitude da aceleração a classificadores treinados \cite{lara2013}.
Este trabalho adota deliberadamente a forma mais simples: um limiar sobre o desvio da magnitude da aceleração em relação a 1\,$g$.
O objetivo não é classificar atividades, e sim distinguir ``em repouso'' de ``em movimento'' de forma explicável e testável.

\subsection{Fadiga de alarmes}

Sendelbach e Funk descrevem a fadiga de alarmes como a dessensibilização da equipe diante de alarmes frequentes e majoritariamente sem ação necessária, com risco à segurança do paciente \cite{sendelbach2013}.
Entre as medidas recomendadas estão ajustar limiares e atrasos aos pacientes e às situações.
Os mecanismos da regra proposta seguem essa direção: persistência para não alertar por um pico isolado, histerese para não oscilar perto do limiar e supressão por contexto de atividade.

\subsection{Mensageria e processamento de eventos}

O padrão produtor-consumidor com um canal intermediário desacopla quem gera mensagens de quem as processa e permite absorver rajadas \cite{hohpe2003}.
Com \emph{um único consumidor}, o canal também impõe uma ordem total de processamento, o que simplifica a regra com estado: não há duas avaliações concorrentes sobre a mesma janela.
Brokers como MQTT foram usados no Marco 2 do projeto.
No Marco 3, a comunicação passou para HTTP, para deixar explícitos os códigos de autenticação (401) e autorização (403), e a fila virou uma estrutura em memória dentro da API.

\subsection{Segurança em IoT}

O HMAC \cite{rfc2104} produz um código de autenticação de mensagem a partir de uma chave compartilhada, que garante integridade e autenticidade sem transmitir a chave.
Sozinho, porém, não garante confidencialidade, e um código antigo pode ser repetido se nada ligar a mensagem ao tempo.
O TLS~1.3 \cite{rfc8446} protege o canal e autentica o servidor por certificado.
Tokens do tipo \emph{Bearer} \cite{rfc6750} dão acesso a quem os possui e, por isso, só podem trafegar em canal protegido.
Para armazenar senhas, funções de derivação de chave caras, como o scrypt \cite{rfc7914}, dificultam ataques de dicionário contra um eventual vazamento do hash.

\subsection{Trabalhos relacionados e posicionamento}

As soluções revisadas em \cite{islam2015} costumam concentrar a decisão em um gateway ou na nuvem, que recebe dados de sensores no corpo.
Os trabalhos de reconhecimento de atividade \cite{lara2013} usam o smartphone ou sensores dedicados para \emph{classificar} a atividade como fim.
O protótipo aqui descrito se diferencia em três aspectos, compatíveis com o escopo didático:
\begin{itemize}
  \item usa a atividade como \emph{contexto de supressão} de um alerta fisiológico, e não como resultado;
  \item trata o contexto como um fato com prazo de validade, e não como um estado permanente;
  \item torna verificável cada regra de decisão e de segurança por testes automatizados.
\end{itemize}
O trabalho não tem pretensão clínica: os sinais são sintéticos e os limiares, escolhidos para demonstração.

% ------------------------------------------------------------------
\section{Proposta e arquitetura}\label{sec:arq}

\subsection{Cenário e requisitos}

Um paciente fictício (\texttt{patient-0042}) veste uma roupa com um ESP32 (\texttt{esp32-borda-01}) que mede FC e aceleração uma vez por segundo.
Um cuidador (\texttt{cuidador01}) acompanha o paciente pelo celular ou pelo computador.
A Tabela~\ref{tab:req} resume os requisitos do quarto marco e onde cada um é atendido.

\begin{table}[ht]
\centering\small
\caption{Requisitos do Marco 4 e como o protótipo os atende.}\label{tab:req}
\begin{tabularx}{\textwidth}{@{}p{3.6cm}L@{}}
\toprule
\textbf{Requisito} & \textbf{Atendimento} \\
\midrule
Dois produtores, decisão e resposta & ESP32 (lotes de FC e aceleração) e app Android (movimento); decisão em \texttt{regra.py}; resposta: estado no app e no painel, LED e buzzer, e reconhecimento do alerta com confirmação (ACK) \\
Mudança de contexto & Movimento informado pelo celular nos últimos 6\,s suprime ou encerra o alerta de FC alta \\
Persistência & TimescaleDB com três \emph{hypertables}: leituras, transições de estado e comandos \\
Segurança & HMAC no dispositivo, sessão no painel, token no app; vínculo dispositivo--paciente--cuidador; HTTPS no app \\
Falha e recuperação & Modo local no ESP32, estado \texttt{DADOS\_INSUFICIENTES} na API, app que mantém o último estado e retoma sozinho \\
Aplicativo móvel & App nativo em Kotlin + Jetpack Compose, com login, estado, reconhecimento do alerta e acelerômetro \\
\bottomrule
\end{tabularx}
\end{table}

\subsection{Visão geral}

A Figura~\ref{fig:arq} mostra os componentes e os fluxos.
Há dois ciclos produtor-consumidor.
No \emph{ciclo 1}, o ESP32 e o app produzem leituras, que a API consome por uma única fila FIFO e transforma em estado.
No \emph{ciclo 2}, o cuidador produz um comando de reconhecimento, o ESP32 o consome ao consultar a API, silencia o buzzer e produz uma confirmação (ACK), que a API consome e repassa às interfaces.

\begin{figure}[ht]
\centering
\begin{tikzpicture}[
  font=\footnotesize,
  caixa/.style={draw, rounded corners=3pt, align=center, minimum height=9mm, inner sep=4pt, fill=#1},
  caixa/.default={blue!6},
  seta/.style={-{Stealth[length=2mm]}, thick},
  rot/.style={font=\scriptsize, align=center, fill=white, inner sep=1pt}]
  \node[caixa=orange!10] (esp) {ESP32 (Wokwi)\\FC sintética + MPU6050\\LED, buzzer, modo local};
  \node[caixa=red!8, right=3.2cm of esp] (app) {App Android\\acelerômetro, estado,\\reconhecer alerta};
  \node[caixa=gray!10, below=1.3cm of $(esp)!0.5!(app)$, minimum width=5cm] (tun) {Túnel Cloudflare (URL temporária)};
  \node[caixa, below=0.9cm of tun, minimum width=5cm] (api) {FastAPI: HMAC / token, vínculo de permissão};
  \node[caixa=yellow!15, below=0.7cm of api, minimum width=5cm] (fila) {Fila FIFO \texttt{asyncio.Queue} (128 itens)};
  \node[caixa, below=0.7cm of fila, minimum width=5cm] (regra) {Consumidor único + \texttt{regra.py}};
  \node[caixa=green!10, below left=0.9cm and -1.2cm of regra] (db) {TimescaleDB\\leituras, eventos\_estado,\\comandos};
  \node[caixa=violet!8, below right=0.9cm and -1.2cm of regra] (ui) {Painel web\\(WebSocket)};
  \draw[seta] (esp) -- node[rot, left=2pt, pos=0.45]{lotes HTTP + HMAC\\comando / ACK} (tun);
  \draw[seta] (app) -- node[rot, right=2pt, pos=0.45]{movimento, estado,\\comando: HTTPS + token} (tun);
  \draw[seta] (tun) -- (api);
  \draw[seta] (api) -- (fila);
  \draw[seta] (fila) -- (regra);
  \draw[seta] (regra) -- (db);
  \draw[seta] (regra) -- (ui);
  \draw[seta, dashed] (ui.east) -| ++(0.6,0) |- node[rot, pos=0.25, right]{reconhecer\\alerta} (api.east);
\end{tikzpicture}
\caption{Arquitetura do Marco 4. Os dois produtores chegam pela mesma fila; a resposta volta ao ESP32 pela consulta periódica de comandos.}\label{fig:arq}
\end{figure}

\subsection{Distribuição das responsabilidades}

A distribuição segue o princípio de colocar cada decisão onde está a informação necessária para tomá-la:
\begin{itemize}
  \item \textbf{ESP32.} Mede, atua e mantém uma regra mínima para quando a rede cai: se a API não responde por 8\,s, entra em \emph{modo local} e alerta sozinho com FC acima de 100\,bpm em repouso. A regra completa saiu do firmware no Marco 3, porque em Python ela é testável sem simulador.
  \item \textbf{API.} É o único ponto em que os dois produtores se encontram, por isso concentra a decisão. Um único processo com um único consumidor garante uma ordem total dos eventos.
  \item \textbf{Aplicativo.} Produz o contexto (movimento) e serve de interface do cuidador. Não decide nada: mostra o estado calculado pela API.
  \item \textbf{Banco temporal.} Fica em um serviço gerenciado na nuvem, para que o histórico sobreviva ao reinício da máquina da API.
\end{itemize}

\subsection{Contratos}

O ESP32 envia lotes com até 10 leituras (na prática, 3 por lote, um lote a cada 3\,s).
Cada lote leva o identificador do dispositivo, o do paciente e um \texttt{loteId} para deduplicação.
O app envia uma leitura a cada 2\,s, com a intensidade do movimento e o instante de captura:

\begin{lstlisting}
POST /api/app/leituras          Authorization: Bearer <token>
{"entityId": "patient-0042", "intensidade": 0.42, "capturadoEmMs": 1764252000123}
-> 200 {"status": "recebido", "estado": "NORMAL", "movimento": "ACTIVE"}
-> 422 {"erro": "Leitura descartada", "motivo": "atrasada", "atrasoMs": 15000}
\end{lstlisting}

A intensidade é o desvio da magnitude da aceleração em relação à gravidade, em unidades de $g$:
\begin{equation}
  i = \left|\, \frac{\sqrt{a_x^2 + a_y^2 + a_z^2}}{9{,}81} - 1 \,\right|.
\end{equation}
O app envia o maior valor de $i$ observado desde o envio anterior, para que uma sacudida curta não se perca entre duas amostras.
Os contratos completos estão versionados no repositório, em \texttt{marco4/contrato/}.

% ------------------------------------------------------------------
\section{Implementação}\label{sec:impl}

\subsection{Dispositivo}

O firmware do ESP32 roda no simulador Wokwi.
Nenhum sensor real entrega a FC pronta: um sensor óptico (PPG) entrega a onda de pulso, e o dispositivo conta os batimentos.
Por isso, um chip próprio do Wokwi simula o sensor e gera a onda, com o ritmo definido por um potenciômetro, e o ESP32 lê a onda a 50\,Hz, detecta cada batimento e calcula a FC pela mediana dos últimos cinco intervalos; 3\,s sem batimento tornam a leitura inválida.
Um MPU6050 virtual fornece a aceleração.
Pelo monitor serial, é possível escolher cenários de normalidade, FC alta em repouso, sensor fora do dedo, FC alta com movimento, ciclo automático e falha de rede; os cenários só alteram o coração simulado, e a FC continua medida na onda.
Cada lote é assinado com HMAC-SHA256 sobre a concatenação do horário (sincronizado por NTP), do método, da rota e do corpo da requisição.
A cada 3\,s, o dispositivo também consulta se há um comando pendente; ao recebê-lo, silencia o buzzer e envia o ACK.

\subsection{API e fila}

A API é escrita em Python com FastAPI e roda em um único processo Uvicorn.
Cada requisição de produtor passa por três etapas antes de entrar na fila:
\begin{itemize}
  \item autenticação: assinatura do dispositivo ou token do app;
  \item autorização: o dispositivo ou o token pertence ao paciente informado;
  \item validação do conteúdo: tipos, faixas e tamanho do lote.
\end{itemize}
Os itens válidos entram em uma \texttt{asyncio.Queue} de capacidade 128, rotulados com a origem (\texttt{esp32} ou \texttt{celular}).
Um único consumidor retira os itens em ordem FIFO, aplica a regra e, se houver banco configurado, grava o lote e a eventual transição.
Por fim, publica o novo estado aos navegadores conectados por WebSocket \cite{rfc6455}.
A resposta HTTP aguarda o processamento, e por isso o produtor recebe o estado calculado na própria resposta.
Com a fila cheia, a API responde 503.

\subsection{Regra de decisão}

A regra mantém uma janela das últimas 10 amostras de FC e das últimas 5 de movimento.
A Tabela~\ref{tab:param} lista os parâmetros, reduzidos para que cada transição apareça em uma demonstração de poucos minutos.
A Figura~\ref{fig:estados} mostra os estados e as transições.

\begin{table}[ht]
\centering\small
\caption{Parâmetros da regra.}\label{tab:param}
\begin{tabularx}{\textwidth}{@{}lrL@{}}
\toprule
\textbf{Parâmetro} & \textbf{Valor} & \textbf{Função} \\
\midrule
Janela de FC / de movimento & 10 / 5 amostras & Média móvel e movimento predominante \\
Cobertura mínima & 70\% & Abaixo disso: \texttt{DADOS\_INSUFICIENTES} \\
Limiar de entrada / saída & 100 / 92 bpm & Histerese \\
Persistência & 2 avaliações & Evita alerta por um lote isolado \\
Tempo mínimo por estado & 4 s & Evita oscilação \\
Silêncio do dispositivo & 8--30 s & 2,5\,$\times$ mediana dos intervalos; sem lote: \texttt{DADOS\_INSUFICIENTES} \\
Repouso & $i \le 0{,}15\,g$ & Mesmo limiar para MPU6050 e celular \\
Validade do contexto do celular & 6 s & Depois disso, vale só o ESP32 \\
Atraso máximo / adiantamento & 10 s / 5 s & Leitura do celular fora disso: 422 \\
\bottomrule
\end{tabularx}
\end{table}

\begin{figure}[ht]
\centering
\begin{tikzpicture}[font=\footnotesize, node distance=2.2cm and 4.2cm,
  estado/.style={draw, rounded corners=4pt, minimum width=3.4cm, minimum height=9mm, align=center},
  seta/.style={-{Stealth[length=2mm]}, thick}, rot/.style={font=\scriptsize, align=center}]
  \node[estado, fill=gray!12] (di) {DADOS\_INSUFICIENTES};
  \node[estado, fill=green!12, below left=2.4cm and 0.4cm of di] (n) {NORMAL};
  \node[estado, fill=red!12, below right=2.4cm and 0.4cm of di] (a) {ATENCAO};
  \draw[seta] (di.south west) -- node[rot, right=3pt, pos=0.4]{cobertura\\$\ge$ 70\%} (n.north east);
  \draw[seta] ([yshift=2mm]n.east) -- node[rot, above]{média $>$ 100 bpm, em repouso,\\2 lotes seguidos, $\ge$ 4 s} ([yshift=2mm]a.west);
  \draw[seta] ([yshift=-2mm]a.west) -- node[rot, below]{média $\le$ 92 bpm \textbf{ou} em movimento\\(ESP32 ou celular), $\ge$ 4 s} ([yshift=-2mm]n.west -| n.east);
  \draw[seta] (a) -- node[rot, right=4pt]{cobertura $<$ 70\%\\ou 8 s sem lote} (di);
  \draw[seta] ([xshift=-8mm]n.north) to[bend left=35] node[rot, left=3pt]{8 s sem lote} (di.west);
\end{tikzpicture}
\caption{Estados do paciente. A transição de ATENCAO para NORMAL causada pelo celular é registrada com a razão \texttt{movimento\_no\_celular}.}\label{fig:estados}
\end{figure}

O contexto do celular entra no cálculo do movimento da seguinte forma.
Seja $m_E$ o movimento predominante nas 5 últimas amostras do ESP32 e $(m_C, t_C)$ a última leitura aceita do celular, com seu instante.
O movimento usado pela regra no instante $t$ é
\begin{equation}
  m(t) = \begin{cases}
    \texttt{ACTIVE}, & \text{se } m_C = \texttt{ACTIVE} \text{ e } t - t_C \le 6\,\text{s},\\
    m_E, & \text{caso contrário.}
  \end{cases}
\end{equation}
Só o movimento \emph{ativo} do celular sobrepõe o do ESP32.
Assim, um celular parado sobre a mesa não anula o movimento detectado na roupa.

Duas restrições completam a combinação.
A primeira: uma leitura do celular pode \emph{suprimir} ou \emph{encerrar} um alerta, mas não conta para a persistência de entrada, de modo que o alerta continua exigindo 2 lotes do ESP32.
A segunda: com o ESP32 em silêncio, a leitura do celular não reavalia a janela antiga, para não tirar o paciente de \texttt{DADOS\_INSUFICIENTES} com dados velhos.
As duas restrições foram acrescentadas depois de a avaliação revelar o defeito discutido na Seção~\ref{sec:result}.

\subsection{Validade e atraso das leituras do celular}

A rede móvel pode atrasar ou reordenar requisições.
Uma leitura que chega tarde descreve um movimento que já passou e não deve mudar o contexto atual.
Por isso, a API compara o instante de captura informado pelo app com o próprio relógio e descarta, com 422, leituras de três tipos:
\begin{itemize}
  \item atrasadas em mais de 10\,s;
  \item adiantadas em mais de 5\,s, o que indica relógio do celular errado;
  \item mais antigas que a última leitura aceita.
\end{itemize}
Leituras descartadas não entram na fila e são contadas no estado, para diagnóstico.

\subsection{Persistência}

Quando a variável \texttt{DATABASE\_URL} está configurada, um adaptador grava os dados em três tabelas convertidas em \emph{hypertables} do TimescaleDB \cite{timescale}, particionadas pelo tempo:
\begin{itemize}
  \item \texttt{leituras}: cada amostra do ESP32;
  \item \texttt{eventos\_estado}: cada transição de estado, com a razão, a média e a cobertura;
  \item \texttt{comandos}: cada mudança de status de um comando (\texttt{PENDENTE}, \texttt{ENTREGUE}, \texttt{EXECUTADO}).
\end{itemize}
As gravações rodam em uma \emph{thread} auxiliar, para não bloquear o laço de eventos, e uma falha no banco não impede o alerta: a API continua decidindo em memória e contabiliza as gravações perdidas.
As leituras do celular não são gravadas, apenas as transições que elas causam.
O valor do movimento perde sentido depois dos 6\,s de validade, e a transição com a razão \texttt{movimento\_no\_celular} já registra o efeito do contexto.

\subsection{Segurança}

A Tabela~\ref{tab:seg} resume como cada participante é identificado e o que pode fazer.
Os segredos ficam fora do Git: a chave do dispositivo, o segredo do servidor e o hash scrypt \cite{rfc7914} da senha do cuidador são gerados localmente por um script, em arquivos ignorados pelo controle de versão.
Uma chave de demonstração chegou a ser versionada no Marco 3, foi removida do repositório e substituída.

\begin{table}[ht]
\centering\small
\caption{Identidade, canal e permissões de cada participante.}\label{tab:seg}
\begin{tabularx}{\textwidth}{@{}p{2.3cm}LLL@{}}
\toprule
\textbf{Participante} & \textbf{Identidade} & \textbf{Canal} & \textbf{Permissões} \\
\midrule
ESP32 & HMAC-SHA256 com chave própria e horário (janela de 120\,s) & HTTP até a borda do túnel (limitação declarada) & Publicar leituras e buscar ou confirmar comandos só do seu paciente \\
Painel web & Senha (scrypt) e cookie de sessão assinado, \texttt{SameSite=Lax} & HTTP local ou HTTPS pelo túnel & Ver o seu paciente e criar comandos \\
App Android & Senha e token \emph{Bearer} assinado, válido por 8\,h \cite{rfc6750} & Só HTTPS, com certificado validado pelo Android & Ver o seu paciente, criar comandos e enviar movimento \\
\bottomrule
\end{tabularx}
\end{table}

O token do app é assinado com o segredo do servidor e carrega o usuário e o paciente.
Por isso, o servidor não guarda sessões do app e nenhuma credencial administrativa vai no APK.
A contrapartida é que um token não pode ser revogado individualmente antes de expirar.
O manifesto do app proíbe tráfego em texto claro, e a conexão usa as autoridades certificadoras do sistema, o que valida a identidade do serviço \cite{rfc8446}.
No ESP32, o \emph{handshake} TLS não funcionou no sketch do grupo, e a comunicação ficou em HTTP com HMAC.
Esse trecho tem integridade e autenticidade, mas não sigilo, e admite repetição de mensagens dentro da janela de 120\,s, atenuada pela deduplicação por \texttt{loteId}.

\subsection{Aplicativo Android}

O app é nativo, escrito em Kotlin com Jetpack Compose, e não usa bibliotecas de rede além das do próprio Android.
Tem duas telas:
\begin{itemize}
  \item \textbf{Login.} Recebe a URL HTTPS do túnel, o usuário e a senha e obtém o token.
  \item \textbf{Estado.} Consulta a API a cada 2\,s e trata as três situações exigidas: carregamento, falha de comunicação (mantendo o último estado na tela) e resultado. Oferece o botão \emph{Reconhecer alerta}, habilitado só em \texttt{ATENCAO}.
\end{itemize}
Em paralelo, o \texttt{SensorManager} entrega as amostras do acelerômetro \cite{androidsensors}, e o app envia a intensidade a cada 2\,s.
Uma resposta 401 leva de volta ao login.

% ------------------------------------------------------------------
\section{Metodologia de avaliação}\label{sec:metodo}

A avaliação responde a quatro perguntas, alinhadas às verificações pedidas para o Marco 4:
\begin{description}[leftmargin=0.6cm, style=nextline]
  \item[Q1 -- Fluxo e contexto] Os dois produtores chegam à decisão e à resposta, e o contexto do celular altera a decisão como projetado?
  \item[Q2 -- Persistência] Leituras, transições e comandos são gravados e recuperados?
  \item[Q3 -- Segurança] Identidades inválidas e operações fora da permissão são rejeitadas?
  \item[Q4 -- Falha e recuperação] O sistema reconhece a falta de dados e volta ao normal quando eles retornam?
\end{description}

Foram usados quatro instrumentos.

\textbf{Testes automatizados.}
São 49 testes em Python (\texttt{unittest}, executados com \texttt{pytest}) que exercitam a regra pura e a API completa dentro do processo, com credenciais geradas por teste, relógio monotônico controlado e o banco substituído por um objeto simulado.
No app, outros 15 testes de unidade em Kotlin (JUnit, executados na JVM, sem emulador) cobrem a lógica que não depende do Android: a intensidade do movimento, a exigência de HTTPS, o corpo da leitura e a interpretação da resposta de estado.
A integração contínua (GitHub Actions) executa as duas suítes a cada envio à \texttt{main}, com Python 3.10 e 3.12 para a API, gera o APK e guarda os relatórios JUnit como evidência.

\textbf{Medições reprodutíveis.}
O script \texttt{avaliar.py} mede quatro aspectos:
\begin{enumerate}[label=(\alph*)]
  \item a latência de processamento dentro da API, com 300 requisições de cada tipo e sem rede;
  \item o custo do login;
  \item os tempos de decisão da regra, com relógio simulado em que o ESP32 envia um lote a cada 3\,s e o celular uma leitura a cada 2\,s;
  \item a taxa de descarte de 2000 leituras sintéticas do celular, com atraso de rede exponencial de média 2\,s e 5\% de leituras retidas entre 12 e 30\,s, com semente fixa.
\end{enumerate}
O ambiente de medição foi um notebook com Windows 11 e Python 3.14.

\textbf{Ensaio do aplicativo no emulador.}
O APK foi instalado no emulador do Android Studio (API 37) e conectado à API pelo túnel HTTPS real.
Um script fez o papel do ESP32, enviando lotes assinados com FC de 120\,bpm a cada 3\,s, buscando o comando e devolvendo o ACK.
O movimento foi injetado no acelerômetro do emulador, e a rede, desligada e religada pelo próprio sistema.

\textbf{Ensaio com o sistema completo.}
O ensaio usa o TimescaleDB na nuvem, o túnel, o Wokwi e o APK instalado em um celular.
Seu roteiro, com entrada, resultado esperado e campo para o resultado observado, está no documento de verificação do repositório.
Ele cobre o que os testes não alcançam: a interface do app, o certificado real do túnel e a preservação dos dados no banco depois de reiniciar a API.

% ------------------------------------------------------------------
\section{Resultados e discussão}\label{sec:result}

\subsection{Testes automatizados}

Os 49 testes da API e os 15 do app passaram localmente e no CI.
A Tabela~\ref{tab:testes} relaciona os módulos de teste às perguntas de avaliação.

\begin{table}[ht]
\centering\small
\caption{Testes automatizados por módulo e pergunta de avaliação.}\label{tab:testes}
\begin{tabularx}{\textwidth}{@{}lrL@{}}
\toprule
\textbf{Arquivo} & \textbf{Testes} & \textbf{O que comprova} \\
\midrule
\texttt{test\_regra.py} & 12 & Q1, Q4: janela, cobertura, persistência, histerese, movimento, silêncio \\
\texttt{test\_api.py} & 12 & Q1, Q3, Q4: ciclos 1 e 2; HMAC inválido, corpo alterado, mensagem antiga (401); dispositivo ou paciente alheio (403); papéis \\
\texttt{test\_melhorias.py} & 6 & Q1: ordem FIFO e contadores; WebSocket só com sessão; diagnóstico sem dados pessoais \\
\texttt{test\_consolidacao.py} & 4 & Q2: gravação de leituras, transições e dos três status de comando; consulta ao histórico; segredos gerados fora do Git \\
\texttt{test\_app\_mobile.py} & 15 & Q1, Q3, Q4: token inválido, forjado ou expirado (401); outro paciente (403); atraso (422); supressão e encerramento por contexto; expiração; regressões da Seção 6.3 \\
\texttt{LogicaTest.kt} (app) & 15 & Q1, Q3: intensidade do acelerômetro; recusa de URL sem HTTPS; corpo conforme o contrato; campos nulos e botão só em \texttt{ATENCAO} \\
\bottomrule
\end{tabularx}
\end{table}

\subsection{Latência de processamento}

A Tabela~\ref{tab:lat} mostra o tempo entre a chegada da requisição e a resposta, medido dentro do processo.
Esse tempo inclui a verificação da credencial, a passagem pela fila, a regra e a serialização do estado.
Ficam de fora a rede, o túnel e o banco.

\begin{table}[ht]
\centering\small
\caption{Latência dentro da API (300 requisições; login: 20).}\label{tab:lat}
\begin{tabular}{@{}lrrr@{}}
\toprule
\textbf{Operação} & \textbf{Média (ms)} & \textbf{Mediana (ms)} & \textbf{p95 (ms)} \\
\midrule
Lote do ESP32 (HMAC + fila + regra) & 5,17 & 4,45 & 8,09 \\
Leitura do app (token + fila + regra) & 4,19 & 4,08 & 5,00 \\
Consulta de estado pelo app & 5,07 & 4,76 & 5,97 \\
Login do app (scrypt, $N=2^{14}$) & 65,67 & 64,84 & 72,68 \\
\bottomrule
\end{tabular}
\end{table}

O processamento leva poucos milissegundos, três ordens de grandeza abaixo dos intervalos de envio (2 a 3\,s).
Com um único consumidor, uma estimativa grosseira da capacidade é de cerca de 200 itens por segundo, muito acima das menos de 1 mensagem por segundo de um paciente.
O login é cerca de 15 vezes mais caro que as demais operações, o que é esperado: o custo do scrypt é intencional e incide uma vez a cada 8\,h por aparelho.
Em uso real, a latência percebida será dominada pela rede e pelo túnel, que não entram nesta medição.

\subsection{Tempos de decisão e efeito do contexto}

A Tabela~\ref{tab:regra} resume o comportamento da regra com relógio simulado.

\begin{table}[ht]
\centering\small
\caption{Comportamento da regra com relógio simulado (lote a cada 3\,s, celular a cada 2\,s).}\label{tab:regra}
\begin{tabularx}{\textwidth}{@{}Lr@{}}
\toprule
\textbf{Situação} & \textbf{Resultado} \\
\midrule
Início com FC normal até sair de \texttt{DADOS\_INSUFICIENTES} (janela cheia) & 9\,s \\
FC sobe de 80 para 120\,bpm em repouso, até \texttt{ATENCAO} & 9\,s \\
FC cai de 120 para 95\,bpm (entre os limiares) & alerta mantido \\
FC cai de 95 para 85\,bpm, até \texttt{NORMAL} & 2\,s \\
Alertas em 10\,min de FC 120\,bpm, sem o celular & 1 \\
Alertas em 10\,min de FC 120\,bpm, com o celular em movimento & 0 \\
Alerta ativo e o celular começa a balançar, até \texttt{NORMAL} & imediato (na leitura) \\
Celular volta a ficar parado, até o alerta voltar & 3\,s \\
Celular para de enviar, até o alerta voltar & 9\,s \\
\bottomrule
\end{tabularx}
\end{table}

O alerta surge 9\,s depois da subida da FC, tempo para a média da janela passar de 100\,bpm e a condição persistir por dois lotes.
A histerese funciona como projetado: 95\,bpm, valor entre a saída (92) e a entrada (100), mantém o alerta.
O efeito do contexto é o pretendido: com o celular indicando movimento, nenhum alerta é gerado em 10 minutos de FC alta.
Quando o alerta já está ativo, a primeira leitura de movimento o encerra, porque o tempo mínimo no estado já passou.
Os dois últimos casos mostram a validade do contexto em ação.
Se o celular passa a informar repouso, o contexto termina na hora e o alerta volta no próximo lote.
Se o celular simplesmente para de enviar, o contexto expira em 6\,s e o alerta volta depois de dois lotes, totalizando 9\,s.

\paragraph{Defeito encontrado pela avaliação.}
Na primeira execução da medição, o alerta voltava em \emph{0\,s} quando o celular parava.
A causa: as leituras do celular, a cada 2\,s, também reavaliavam a regra e contavam para a persistência de duas avaliações, que fora projetada para contar lotes do ESP32.
Com o segundo produtor, a persistência passou a ser atingida por duas mensagens no mesmo segundo.
A mesma análise mostrou outro problema: uma leitura do celular durante o silêncio do ESP32 poderia reavaliar a janela antiga e tirar o paciente de \texttt{DADOS\_INSUFICIENTES}.
As duas restrições da Seção~4.3 corrigiram esses defeitos, e dois testes de regressão, que falham sem a correção, foram incluídos na suíte.
O episódio ilustra um risco geral de sistemas com múltiplos produtores: um parâmetro definido em ``número de avaliações'' muda de significado quando um novo produtor passa a disparar avaliações.
A lição prática é definir persistência e validade em termos da \emph{fonte} ou do \emph{tempo}, nunca do número de mensagens.

\subsection{Validade das leituras}

Das 2000 leituras sintéticas, 1564 (78,2\%) foram aceitas, 114 (5,7\%) descartadas por atraso acima de 10\,s e 322 (16,1\%) descartadas por chegarem fora de ordem.
O descarte por atraso corresponde quase exatamente às leituras retidas de propósito.
O descarte por ordem é alto porque o atraso médio (2\,s) é igual ao intervalo entre leituras, e qualquer leitura ultrapassada por uma mais nova é descartada.
Para o contexto, essa política é adequada: a leitura mais nova já descreve o movimento atual, e a mais velha não acrescenta informação.
Se o mesmo critério fosse aplicado a medidas cujo histórico importa, como as próprias amostras de FC, a perda seria inaceitável, e seria preciso reordenar em vez de descartar.

\subsection{Persistência, segurança e falha}

Os testes de persistência comprovam que leituras, transições e os três status de cada comando são gravados e que o histórico é consultado pela API.
Como o banco é simulado nos testes, a comprovação com o TimescaleDB real fica com o ensaio: contar as linhas, reiniciar a API e contar de novo.
Nos testes de segurança, todas as identidades inválidas recebem 401 e todas as operações fora da permissão recebem 403:
\begin{itemize}
  \item token inválido, forjado ou expirado;
  \item assinatura errada ou corpo alterado;
  \item mensagem antiga;
  \item outro paciente, dispositivo não vinculado, ESP32 criando comando ou cuidador forjando ACK.
\end{itemize}
Nos testes de falha, o silêncio do ESP32 leva a \texttt{DADOS\_INSUFICIENTES}, e o celular não tira o paciente desse estado.
A recuperação com o retorno das leituras e o comportamento do app sem rede são verificados no ensaio.

\subsection{Ensaio do aplicativo}

A Figura~\ref{fig:app} mostra quatro momentos do ensaio no emulador.
Com a FC em 120\,bpm e o ESP32 em repouso, o estado passou a \texttt{ATENCAO} 12\,s depois do primeiro lote, como esperado pela janela, pela persistência e pelo tempo mínimo da regra.
O reconhecimento feito no app chegou ao ESP32 simulado no lote seguinte, que devolveu o ACK, e o cartão mostrou \texttt{EXECUTADO}.
Com o acelerômetro marcando 0,83\,$g$, o estado passou a \texttt{NORMAL}, com o movimento vindo do celular, entre dois lotes; parado o movimento, o alerta voltou.
Sem rede, o app manteve o último estado e indicou a falha; com a rede de volta, retomou as consultas sem novo login.
O app também recusou a URL com \texttt{http://} e mostrou a mensagem de senha incorreta para o 401.

\begin{figure}[ht]
\centering
\newcommand{\tela}[2]{\begin{minipage}[t]{0.235\textwidth}\centering
  \includegraphics[width=\linewidth]{fig/app-#1.png}\\[2pt]{\footnotesize #2}\end{minipage}}
\tela{05-atencao}{(a) FC alta em repouso}\hfill
\tela{06-reconhecido}{(b) alerta reconhecido}\hfill
\tela{07-contexto-movimento}{(c) celular em movimento}\hfill
\tela{09-sem-rede}{(d) sem rede}
\caption{Tela de estado do app no ensaio no emulador. Em (c), o movimento do celular encerra o alerta mesmo com a FC em 120\,bpm.}\label{fig:app}
\end{figure}

O ensaio revelou um defeito de apresentação que os testes da API não alcançavam: a biblioteca \texttt{org.json} devolve o texto \texttt{"null"} para campos nulos, e a tela exibia ``Último contato: null''.
A correção passou a ser coberta por um teste de unidade do app.

\subsection{Ensaio com o ESP32 no Wokwi}

Com o firmware no Wokwi, a detecção de batimentos mediu 77--79\,bpm para o coração simulado a 78\,bpm, 136\,bpm para 135 e 125\,bpm para o potenciômetro no meio da escala.
Com o sensor fora do dedo, a FC ficou inválida em cerca de 3\,s, e a API passou a \texttt{DADOS\_INSUFICIENTES}.
Com o artefato de movimento, a FC medida oscilou entre 97 e 120\,bpm para um coração a 135, o erro típico de sensores ópticos durante o movimento; ali a FC alta não gera alerta, porque o movimento a explica.
O ensaio trouxe dois achados.
O simulador roda a cerca de 40\% da velocidade real, e os lotes de 3\,s chegavam a cada 7--10\,s; com o limite fixo de 8\,s, o estado alternava por silêncio.
O limite passou a acompanhar o ritmo dos lotes (Tabela~\ref{tab:param}), com teto de 30\,s para que uma queda real continue detectada.
Além disso, o túnel passou a redirecionar \texttt{http://} para HTTPS, e o ESP32 passou a falar com a API pelo gateway local do simulador, sem sair da máquina.
Nesse ensaio, o reconhecimento feito no app silenciou o buzzer do ESP32 simulado, e a queda de rede levou ao modo local, com recuperação ao voltar.

\subsection{Limitações}

\begin{itemize}
  \item \textbf{Canal do dispositivo.} O trecho entre o ESP32 e o túnel não tem sigilo, e a janela de 120\,s admite repetição. Com dados reais, isso exigiria TLS no dispositivo.
  \item \textbf{Volatilidade.} Fila, estado corrente e comando pendente ficam em memória. Depois de reiniciar a API, o estado recomeça vazio, embora o histórico esteja no banco.
  \item \textbf{Validação de campo.} As medições de latência e de tempo foram feitas sem rede, e os ensaios usaram o emulador e o Wokwi, sem banco. Os resultados com o celular real e o TimescaleDB entram na versão final deste relatório.
  \item \textbf{Papel do celular.} Na demonstração, o mesmo celular é o app do cuidador e o sensor do paciente. Num cenário real, o sensor ficaria com o paciente, por exemplo em um relógio.
  \item \textbf{Energia.} O app mantém o acelerômetro ativo e faz duas requisições a cada 2\,s enquanto está aberto. O consumo não foi medido.
  \item \textbf{Validade clínica.} Os sinais são sintéticos e os limiares, de demonstração. Nenhuma conclusão clínica pode ser tirada.
\end{itemize}

% ------------------------------------------------------------------
\section{Conclusão}\label{sec:concl}

Este trabalho apresentou um protótipo de monitoramento ubíquo em que o alerta de FC elevada depende do contexto de atividade do paciente, informado por um segundo produtor, o acelerômetro do smartphone.
A arquitetura junta os dois produtores em uma fila FIFO com um consumidor único, aplica uma regra com estado, tempo e validade de contexto, persiste o histórico em um banco temporal e fecha o ciclo com um comando confirmado pelo dispositivo.
Cada participante tem identidade e permissões próprias, sem segredos no repositório nem no aplicativo.

A avaliação mostrou que o processamento é rápido diante dos intervalos de envio e que o contexto elimina os alertas de FC alta durante o movimento, mantendo-os em repouso.
Mostrou também que acrescentar um produtor muda o significado de parâmetros definidos em número de mensagens.
O defeito encontrado e corrigido reforça o valor de medir a regra com relógio simulado, além de testá-la caso a caso.

Os trabalhos futuros, previstos para a entrega final, são:
\begin{itemize}
  \item TLS no dispositivo;
  \item teste de integração com o TimescaleDB real no CI;
  \item testes instrumentados do aplicativo;
  \item validar no ensaio integrado a recuperação do estado e do comando pendente após reinício, incluindo os limites da fila volátil;
  \item registro dos resultados do ensaio completo com o APK.
\end{itemize}

\section*{Declaração de uso de IA}
O assistente Claude Code (modelo Claude Opus 5.5) foi usado no Marco 4 para planejar as tarefas, implementar as rotas do app e a regra de contexto, escrever testes, configurar o CI, gerar o esqueleto do app, redigir a documentação e esboçar este texto.
O grupo revisou o material.
A validação foi feita pelos testes automatizados, pelas medições reprodutíveis e pelo ensaio do sistema completo.
O defeito da Seção~6.3 foi identificado a partir dessas medições.
Durante o Marco 4, ChatGPT (GPT-6) também foi utilizado para ampliar o painel Android nativo, integrar a consulta ao histórico, implementar a recuperação do snapshot persistido, preparar testes e atualizar a integração contínua. O grupo deve revisar essas alterações e validar os resultados antes da versão final.
O grupo é responsável por compreender, testar e explicar o conteúdo entregue.

\begin{thebibliography}{99}\small
\bibitem{weiser1991} WEISER, M. The computer for the 21st century. \emph{Scientific American}, v.~265, n.~3, p.~94--104, 1991.
\bibitem{satya2001} SATYANARAYANAN, M. Pervasive computing: vision and challenges. \emph{IEEE Personal Communications}, v.~8, n.~4, p.~10--17, 2001.
\bibitem{dey2001} DEY, A.~K. Understanding and using context. \emph{Personal and Ubiquitous Computing}, v.~5, n.~1, p.~4--7, 2001.
\bibitem{abowd1999} ABOWD, G.~D.; DEY, A.~K.; BROWN, P.~J.; DAVIES, N.; SMITH, M.; STEGGLES, P. Towards a better understanding of context and context-awareness. In: \emph{Handheld and Ubiquitous Computing (HUC'99)}. Berlin: Springer, 1999. (LNCS, v.~1707), p.~304--307.
\bibitem{islam2015} ISLAM, S.~M.~R.; KWAK, D.; KABIR, M.~H.; HOSSAIN, M.; KWAK, K.-S. The Internet of Things for health care: a comprehensive survey. \emph{IEEE Access}, v.~3, p.~678--708, 2015.
\bibitem{lara2013} LARA, O.~D.; LABRADOR, M.~A. A survey on human activity recognition using wearable sensors. \emph{IEEE Communications Surveys \& Tutorials}, v.~15, n.~3, p.~1192--1209, 2013.
\bibitem{sendelbach2013} SENDELBACH, S.; FUNK, M. Alarm fatigue: a patient safety concern. \emph{AACN Advanced Critical Care}, v.~24, n.~4, p.~378--386, 2013.
\bibitem{hohpe2003} HOHPE, G.; WOOLF, B. \emph{Enterprise Integration Patterns: designing, building, and deploying messaging solutions}. Boston: Addison-Wesley, 2003.
\bibitem{rfc2104} KRAWCZYK, H.; BELLARE, M.; CANETTI, R. \emph{HMAC: keyed-hashing for message authentication}. RFC 2104. IETF, 1997.
\bibitem{rfc8446} RESCORLA, E. \emph{The Transport Layer Security (TLS) protocol version 1.3}. RFC 8446. IETF, 2018.
\bibitem{rfc6750} JONES, M.; HARDT, D. \emph{The OAuth 2.0 authorization framework: bearer token usage}. RFC 6750. IETF, 2012.
\bibitem{rfc7914} PERCIVAL, C.; JOSEFSSON, S. \emph{The scrypt password-based key derivation function}. RFC 7914. IETF, 2016.
\bibitem{rfc6455} FETTE, I.; MELNIKOV, A. \emph{The WebSocket protocol}. RFC 6455. IETF, 2011.
\bibitem{timescale} TIGER DATA. \emph{TimescaleDB documentation}. Disponível em: \url{https://www.tigerdata.com/docs}. Acesso em: 27 nov. 2026.
\bibitem{androidsensors} ANDROID DEVELOPERS. \emph{Motion sensors}. Disponível em: \url{https://developer.android.com/develop/sensors-and-location/sensors/sensors_motion}. Acesso em: 27 nov. 2026.
\end{thebibliography}

\end{document}
